% Ad Quintum Fratrem 3.2 --- headnote
% ad-quintum-fratrem-03-02, 11 October 54 BC, Rome

\textit{Cicero to Quintus, written from Rome before light on
the fifth day before the Ides of October (11 October 54 BC).
The letter is the close-up of two political moments: the
divinatio that opened the prosecution of Gabinius, and the
canvassing-charges hurricane that descended on all four
candidates for the consulship of 53 BC.}

\textit{Gabinius, returned from his Syrian governorship to
the universal hatred reported in \textit{Q. fr.} 3.1, has been
rolled up in the Senate. The narrative at \textsection 2 is
the only first-person account of an unforgettable scene:
Gabinius, having boasted that he would demand a triumph,
slipping into Rome by night ``as into a city full of
enemies''; then, on the tenth day on which he had to report
the campaign returns, creeping into the Senate at lowest
attendance. With the publicans brought in, with Cicero
slicing him to pieces, Gabinius lost his nerve and called
Cicero ``an exile''; the whole Senate rose to a man, ``with a
shout, as if it would go up against his very body.'' Cicero's
recorded reaction is the small candid moment of the letter:
``nothing more honourable has ever happened to us.'' The
following morning the divinatio against Gabinius would name
the prosecutor for the \textit{maiestas} charge, contested
between Memmius, Ti. Nero, and the two Antonii.}

\textit{Cicero himself does not prosecute Gabinius (\textit{Pro
Rabirio Postumo} of the next year would force him into the
opposite role, defending one of Gabinius's collaborators).
Three reasons are given at \textsection 2: he does not wish
to fight Pompey (Pompey was vehemently working for Gabinius's
acquittal); he has no faith in the jurors of these courts;
and he fears that, with him as accuser, ``something might
happen'' --- the same Pompey-Gabinius backstop that, in fact,
would secure Gabinius's acquittal at the \textit{maiestas}
trial later in October before being undone by his conviction
on extortion. The Greek tag \textit{apoteugma} ---
``failure,'' the technical word from athletic and forensic
contests for the missed throw --- is the candid private
reckoning.}

\textit{The four-way bribery hurricane of \textsection 3
(every candidate prosecuting his rival for ambitus) is the
prelude to the famous interregnum of late 54 BC: with the
courts gummed up and the elections pushed past 1 January, the
state spent the first half of 53 BC without consuls. The
``Lentulus'' of the closing flourish is P. Lentulus Spinther
the elder (Cicero's correspondent of Fam.\ 1), who as
\textit{interrex} of the year prosecuted Gabinius for
\textit{maiestas}; ``Appius'' is Ap. Claudius Pulcher, the
brother of Clodius, who would step into the consulship of 53
without the curiate law that would normally have been required.
The closing domestic note --- the contractors at Quintus's
house ``not handling it undiligently'' --- is the fixed
counter-rhythm of the political letters of 54.}
